\chapter{绪论}
\label{cha:intro}

\section{研究背景与意义}

脑机接口能够在不依赖外周神经和肌肉输出的条件下建立人脑与外部设备之间的信息通道，在神经康复、辅助交互、智能可穿戴和新型人机协同等方向具有重要应用价值 \cite{chaudhary2016brain, lotte2018review}。其中，基于脑电（Electroencephalography, EEG）的非侵入式脑机接口因采集成本较低、使用风险较小、便于重复实验和系统集成，成为当前研究与应用开发最活跃的技术路线之一。对本文而言，这一点尤其重要，因为一旦脑机接口希望进入辅助交互与便携式使用场景，安全性、可重复性和系统可集成性往往比单次实验中的极限性能更具现实意义。

从任务范式看，脑电脑机接口已经形成了较为清晰的技术谱系。以 P300、稳态视觉诱发电位（Steady-State Visual Evoked Potential, SSVEP）和运动想象（Motor Imagery, MI）为代表的范式，分别对应事件相关电位检测、持续视觉刺激驱动识别和主动意图解码等不同交互路径 \cite{farwell1988talking, bin2011high, pfurtscheller1999event}。其中，运动想象无需真实肢体动作即可诱发可识别的感觉运动节律变化，也不依赖持续外部刺激维持，因此长期被视为主动控制类脑机接口中的典型范式；但与此同时，它对个体训练、状态稳定性和解码策略的依赖也更为突出。

无论采用何种范式，脑电解码都面临若干共性难题。首先，头皮脑电信号幅值较弱，易受肌电、眼电和环境噪声污染，整体信噪比较低 \cite{goncharova2003emg, lotte2018review}。其次，脑电信号具有显著的非平稳性，跨被试差异、跨 session 漂移以及试次质量波动都会直接影响模型泛化。再次，脑机接口评价通常不仅涉及静态离线准确率，还受到标注方式、校准成本、交互延迟和在线反馈机制等因素影响。因此，后续关于模型规模、导联数量、在线适配和系统部署的讨论，本质上都在回应这些共性问题。

近年来，围绕上述问题，脑电解码研究一方面沿着更深的卷积网络、更复杂的自注意力结构以及更大规模训练数据持续推进，另一方面也出现了面向跨被试迁移和表示复用的大规模预训练模型。相关工作试图通过大规模跨数据集学习缓解异质性、样本不足和泛化困难 \cite{Jia24, Wan24b}。但现有证据也表明，这类路线的主要优势仍集中于多导联、离线基准和较重适配条件下的表示迁移，其收益并不能自动等同于低成本、少导联、实时交互和可部署系统中的直接可用性 \cite{Liu26b}。因此，本文并不将“大模型”作为否定对象，而是将其视为当前脑电研究的重要前沿；与此同时，本文更关注在具体系统约束下如何形成够用的模型、合理的任务收束和可落地的实现链路。

随着辅助交互终端、可穿戴设备和日常化人机协同需求的发展，无创脑机接口系统正在逐步突破传统高密度有线实验平台的使用边界。市场上已经出现了以头环、头带、耳周设备等为代表的轻量化 EEG 产品，学术界也在持续探索耳电、前额阵列、低准备度无线采集和针对特定任务的可穿戴 EEG 系统 \cite{He23, Nis22, Arn20, Deb15, Bio22}。这说明“更轻、更易用、更可持续部署”的 EEG 形态并非凭空设想，而是已经在产品层和研究层同时展开的现实趋势。

但便携化、无线化和边缘化并不意味着简单地把实验室设备做小。面向真实使用场景，系统往往需要同时兼顾穿戴便利性、布设时间、无线传输、持续运行、本地响应以及维护成本等多重要求 \cite{Cas19c, Nis22, He23}。在这种条件下，系统通常难以继续沿用实验室中的高密度有线配置，而需要在佩戴复杂度、部署成本与实时响应之间重新取得平衡。相应地，采集端更可能采用简化导联配置和紧凑硬件实现，后端则必须在有限算力、有限存储和有限功耗预算下完成稳定推理，这也使得“实验室最优观测”与“真实场景可部署观测”之间出现了新的张力。

由此带来的挑战可以概括为三个相互耦合的问题束。其一是采集配置与信号稳定性问题：导联数量减少后，可观测信息边界随之收缩，接触状态、运动伪迹、基线漂移和试次质量波动会变得更加敏感。其二是解码任务与算法问题：高密度标准设置下成立的任务定义与模型结构，未必仍然适合少导联系统；系统关注的重点也不再只是“继续提高离线准确率”，而是如何在约束条件下保持足够稳健和足够可用。其三是部署与在线交互问题：有限算力、有限功耗与在线更新需求共同作用，使得模型结构、推理链路和交互反馈机制必须被统一设计，而不能继续作为彼此割裂的模块分别讨论。

基于上述问题意识，本文面向辅助交互与便携式使用场景下的低成本无创脑机接口需求，以运动想象这一典型脑电交互任务为切入点，研究标准高密度离线结论在少导联系统约束下的可迁移性及其落地路径。论文围绕双导采集硬件实现、端侧友好解码方法设计以及在线联动初步验证三个层面展开工作，尝试回答在少导联采集、有限算力和在线交互约束下，什么任务定义更合理、什么模型更够用，以及什么样的系统链路更适合真实部署。

\section{相关研究现状}
\subsection{便携式脑机接口设备、边缘硬件与低成本脑电采集实现研究现状}

从应用形态看，面向辅助交互、可穿戴终端和居家训练场景的无创脑机接口系统，通常需要同时满足便携化、低成本、无线连接和持续运行等要求。这意味着系统设计不再只是追求更完整的空间观测，而是需要在采集复杂度、硬件资源和交互可用性之间进行权衡。相关综述已经表明，可穿戴与无线 EEG 系统在形态上呈现出明显多样化趋势，头环、耳周设备、前额阵列、医疗型减少导联系统以及多种研究原型正在共同推动 EEG 由传统高密度有线平台向更轻量的部署形态扩展 \cite{Cas19c, Nis22, He23}。

从具体路线看，现有研究大致可以分为两类。第一类工作强调轻量设备和低准备度采集形态本身，例如头环式睡眠监测、耳电长期采集以及面向特定临床任务的减少导联系统 \cite{Arn20, Deb15, Bio22}。这类研究直接说明，在某些任务场景中，更易佩戴、更适合长期使用的 EEG 形态已经具备现实需求和技术基础。第二类工作则强调边缘硬件与紧凑型模型的协同可行性。Wang 等在低功耗边缘计算条件下完成了 MI 脑电模型部署，Schneider 等给出了 8 bit 量化 EEGNet 在 PULP 平台上的时延与能效结果，En\'eriz 等则进一步在 Arduino Nano 级平台上完成了 PTQ 部署验证 \cite{Wan20, Sch20, Ene23}。此外，Wan 等通过通道选择与 8 bit 量化的结合，说明在减少输入通道与压缩模型规模的同时，MI 解码仍可保持可接受表现 \cite{Wan22b}；部分相邻时序任务也开始面向 RK3568、NPU 或异构推理引擎展开优化 \cite{You24, Ism23}，为脑电模型的端侧部署提供了平台层面的参考。

总体来看，这类研究已经说明，低成本前端、边缘算力节点和紧凑型模型的组合在原理上是可行的。但多数工作仍更关注“模型能否运行”或“平台是否具备推理能力”，而对真实采集条件下应采用何种导联配置、何种任务定义以及如何形成完整系统链路，仍缺少统一回答。

\subsection{运动想象脑电解码方法研究现状}

在运动想象脑电解码领域，传统机器学习方法长期构成重要基线。以共空间模式（Common Spatial Pattern, CSP）及其滤波器组扩展方法（Filter Bank CSP, FBCSP）为代表的空间滤波路线，能够突出不同类别在空间方差分布上的差异，并与线性判别分析（LDA）、支持向量机（SVM）等浅层分类器结合取得较高的分类效率 \cite{ang2012filter, lotte2018review}。这类方法实现简单、计算代价较低、可解释性较强，但其效果通常依赖预设频带、人工特征与相对稳定的采集条件。

为减少对手工特征工程的依赖，卷积神经网络随后成为脑电解码的重要技术路线。Schirrmeister 等提出的深浅卷积网络说明，端到端学习方式可以直接从原始脑电信号中提取具有判别性的时空特征 \cite{schirrmeister2017deep}；Lawhern 等提出的 EEGNet 则以紧凑型卷积结构兼顾了较低参数规模与较好的跨范式适应能力 \cite{lawhern2018eegnet}。相比传统特征工程流程，CNN 有效推动了脑电解码向自动化表示学习演进，并成为后续轻量化脑电模型的重要基础。

在 CNN 的基础上，Transformer 及其变体被引入脑电解码领域，以增强对长程时序依赖和跨特征关联的建模能力 \cite{vaswani2017attention}。Song 等提出的 EEG Conformer 通过卷积与自注意力协同提取局部时空特征和全局时间依赖 \cite{songEEG2022}，Xie 等进一步展示了 Transformer 结构在原始脑电分类任务中的应用潜力 \cite{xie_transformer-based_2022}。随后，面向脑电序列的高效 Transformer 设计逐步出现，研究重点从“能否使用自注意力”转向“如何在控制复杂度的同时保留关键建模能力”。例如，Wan 等将通道注意力与窗口化机制结合建模时空耦合信息 \cite{wan2023novel}，Hua 等通过卷积与滑动窗口注意力的协同设计增强了多尺度特征建模能力 \cite{hua2023improved}，Keutayeva 等则进一步强调了紧凑型结构在跨被试任务中的部署潜力 \cite{keu2024compact}。

与此同时，近年来还出现了面向脑电大规模预训练与基础模型的研究路线，试图通过跨数据集学习缓解跨被试差异、标签不足和表示迁移困难。相关工作在多任务离线基准上展示了较强的表征复用潜力 \cite{Jia24, Wan24b}，但最新综述也指出，这类模型的主要证据仍集中于多导联、离线和较重适配条件下的表现，其结论尚不能直接等同于少导联、可穿戴和交互式系统中的可部署收益 \cite{Liu26b}。因此，大规模预训练为脑电解码提供了新的重要方向，但其与低成本便携式系统之间仍存在明显的场景差异。

总体来看，运动想象脑电解码方法已经形成了从传统空间滤波、卷积网络到自注意力建模的清晰演进路径，并在标准数据集上不断提升分类性能。但这一路线的主要验证场景仍集中于离线、标准协议和相对高密度导联条件，其结论如何映射到真实系统约束之下，仍有待进一步研究。

\subsection{在线脑机接口、校准机制与交互应用研究现状}

从应用形态看，脑机接口系统并不仅仅由解码模型构成，还同时涉及信号采集、无线传输、提示与反馈、在线校准和用户适应等多个环节。早期运动想象在线研究已经表明，反馈与训练会显著影响脑电模式的可分辨性。Pfurtscheller 等在早期左右手在线想象实验中验证了基于中央区信号的在线判别可行性 \cite{Pfu97}；Neuper 等进一步指出，反馈调节有助于增强左右手感觉运动节律差异 \cite{Neu99}。这说明，在线系统中的任务表现不仅取决于模型本身，也受被试学习过程和交互机制影响。

随后，更面向应用的在线脑机接口研究开始将导联数量、算法结构和行为表现放在统一框架下考察。Meng 等在多轮在线 MI-BCI 实验中比较了小规模传感器配置和大规模传感器配置，结果表明围绕中央区构建的紧凑通道方案在长期训练中仍可获得具有竞争力的在线表现 \cite{Men18}。这类工作说明，便携化和可用性并不一定要求完整保留高密度导联，但也提示在线系统中的任务设计、用户训练和硬件实现不能与解码算法割裂讨论。

近年来，真实在线系统研究与端侧部署研究虽然都在推进，但两者之间仍存在脱节。一部分工作强调在线反馈或被试适应，却未深入讨论硬件资源与部署路径；另一部分工作则提供了模型压缩与设备运行数据，却缺少基于真实在线交互的任务级验证。因此，在线脑机接口的评价不能只停留在离线平均准确率上，系统链路、校准成本、试次质量波动和反馈机制同样需要纳入统一考量。

\section{现有研究存在的不足}

综合前述研究可以看出，当前运动想象脑电解码虽然在模型设计、局部高效注意力和轻量化探索方面取得了明显进展，但对于辅助交互与便携式使用场景下的低成本真实系统目标而言，仍存在以下四方面不足。

第一，场景驱动的系统约束仍缺少整体化分析。现有研究往往分别讨论低成本采集、边缘部署或在线交互中的某一方面，却较少从辅助交互与便携式使用场景出发，统一分析穿戴便利性、无线通信、持续运行、本地响应以及采集稳定性之间的耦合关系。结果是系统需求与算法目标之间的对应关系仍不够清晰。

第二，现有标准任务与真实少导联系统目标之间仍存在明显错位。公开数据集中的高密度四分类设置为算法研究提供了统一基准，但当导联数显著降低时，可观测信息边界随之收缩，原有任务定义并不一定仍然合理。许多工作默认“模型只要足够强就能迁移到更少通道”，却较少系统回答在双导或少导联条件下究竟应保留什么目标、舍弃什么目标，以及应如何构建相应的定量证据链。

第三，算法设计与端侧部署约束之间的协同仍然不足。现有研究往往分别强调高离线精度、模型压缩或硬件加速中的某一方面，但较少在统一框架下同时考虑输入通道规模、结构表达能力、参数规模、量化友好性和推理延迟等约束 \cite{keu2024compact, wan2023novel, hua2023improved}。尤其是在运动想象脑电场景中，“能在公开数据集上取得较高准确率”并不等同于“能够在少导联端侧系统中稳定运行并保持可用表现”。

第四，在线闭环证据与持续校准机制仍有待进一步明确。便携式脑机接口不仅取决于模型推理性能，还同时受到采集链路噪声、BLE 传输时序、试次间漂移和在线校准代价等因素影响。已有研究说明反馈训练、个体差异和试次质量会直接影响在线表现，但关于如何在低成本设备上通过够用的模型、轻量的在线更新和合理的交互设计获得稳定收益，公开研究仍然有限。这也是本文后续在系统集成与在线初步验证部分进一步关注的问题。

\section{本文主要研究内容与技术路线}
\subsection{主要研究内容}

围绕上述问题，本文面向辅助交互与便携式使用场景下的低成本无创脑机接口系统，开展了以下几方面研究工作。

（1）设计并实现双导脑电采集硬件系统。针对辅助交互与便携式使用场景下的低成本采集需求，本文完成了以双导联采集为核心的硬件系统设计与实现，构建了脑电前端、主控、BLE 回传和上位机交互的基础链路，并通过节律验证、标签链路验证和通信稳定性测试，为后续在线实验提供系统基础。

（2）研究面向低通道端侧场景的改进型自注意力脑电解码方法。针对标准高密度解码模型难以直接迁移到少导联端侧系统的问题，本文设计了兼顾局部时序建模能力与部署友好性的改进型自注意力解码方法，在标准条件下完成性能、复杂度与量化友好性分析，并为后续低通道扩展提供算法基础。

（3）构建从标准高密度设置到双导可落地任务的收束路径。围绕 `22导 / 4分类` 标准设置，本文系统分析了导联压缩下的性能退化边界，并在此基础上将目标进一步收束到 `2导 / 2分类` 的左右手运动想象解码任务，从而使算法目标与双导系统可观测信息边界相匹配。

（4）完成双导端侧系统集成与在线联动初步验证。本文以双导硬件为前端、以上位机和可扩展推理节点为原型载体，构建了预训练初始化、轮次级模型更新与跨轮次判别反馈的在线联动机制，并结合伪在线评估和真实在线试跑，对系统的可实现性进行了初步验证。需要指出的是，当前实验阶段的在线验证载体主要为 Linux PC 原型环境，后续目标部署节点则面向 RK3568 / RKNN 一类嵌入式平台。

\subsection{技术路线与论文组织}

全文的技术路线与组织结构如图\ref{fig:org}所示。
\begin{figure}[htbp]
  \centering
  % \includegraphics[width=0.8\textwidth]{org_structure}
  \caption{论文技术路线与组织结构}
  \label{fig:org}
\end{figure}

本文后续各章安排如下：

第 2 章介绍运动想象脑电解码与端侧系统实现所需的关键技术基础，并在比较可选技术路线的基础上说明本文采用的建模与部署路线。

第 3 章围绕双导脑电采集硬件系统的设计与实现展开，重点说明采集前端、主控、无线通信和基础链路验证结果。

第 4 章研究面向低通道端侧场景的改进型自注意力脑电解码方法，在标准条件下建立离线基线，并进一步分析导联压缩下的性能退化边界，完成从 `22导 / 4分类` 到 `2导 / 2分类` 的任务收束与模型适配。

第 5 章进一步将第 4 章得到的低通道解码主线映射到双导系统中，介绍系统集成、在线实现机制、伪在线评估以及真实在线初步验证结果。

第 6 章对全文工作进行总结，并围绕更强的人机协同、自适应学习和更自然的实验范式给出后续展望。
